\documentclass[11pt]{article}
\usepackage[a4paper, top=18mm, bottom=18mm, left=20mm, right=20mm]{geometry}
\usepackage[T2A]{fontenc}
\usepackage[utf8]{inputenc}
\usepackage[ukrainian]{babel}
\usepackage{graphicx}
\usepackage[export]{adjustbox}
\usepackage{amsmath}
\usepackage{amssymb}
\graphicspath{{/app/Quiz/Scripts/2023/ProgAll/15BinTree/}{/app/Quiz/Scripts/2021/Mod2021_3/BinTree/}{/app/Quiz/Scripts/2020/Mod2020_4/BinTree/}}
\setlength{\parindent}{0pt}
\setlength{\parskip}{6pt}
\pagestyle{empty}
\begin{document}
%begin
{\large\textbf{Контрольна робота}}

\textbf{Варіант \No\,9}\hfill \textit{ПІБ:}\,\underline{\hspace{60mm}}
\vspace{4mm}

%beginitem
1. Що буде виведено за виконання фрагменту коду?

%code
print('R', end=' ')
j=1
while j<=9:
    j+=2
print(j)
%code

\vspace{5mm}
%enditem

%beginitem
2. Що буде виведено за виконання фрагменту коду?

%code
j=1
while j<=9:
    j+=2
print(j)
%code

\vspace{5mm}
%enditem

%beginitem
3. %goodpath:Scripts/2018/Mod2018_1/03-good.txt
Відмітити все, що є коректним оператором С++:
( 1 ) return x<=2;

( 2 ) double n.m;

( 3 ) cin>>x+3;

( 4 ) if (1>2) x=10;
\vspace{5mm}
%enditem

%beginitem
4. Що буде виведено за виконання фрагменту коду?
code
__b = 0

class B:
    __b = 1
    
    def __init__(self):
        self.__b = 2
        __b = 3
        B.__b = 4

obj = B()
try:
    print(__b, end=' ')
    print(obj.__b, end=' ')
    print(B.__b, end=' ')
    print(obj.__b, end=' ')
    print(obj._B__b, end=' ')
except:
    print('ERR')
code
\vspace{5mm}
%enditem

%beginitem
5. %goodpath:Scripts/2019/Mod2019_1/Lex/03-good.txt
Відмітити все, що є коректним оператором С++:
( 1 ) return 'a'<'c'

( 2 ) print(sep=":","qwe")

( 3 ) x=*2

( 4 ) float(2)/87
\vspace{5mm}
%enditem

%beginitem
6. Що буде виведено за виконання фрагменту коду?

%code
print('R', end=' ')
b=['0','1','2','3','4']
b[-8:-8]='12'
print(b)
%code

\vspace{5mm}
%enditem

%beginitem
7. %goodpath:Scripts/2019/Mod2019_1/Lex/01-good.txt
Відмітити все, що є однією правильною лексемою:
( 1 ) 'c'

( 2 ) 0123

( 3 ) 0b23

( 4 ) xy
\vspace{5mm}
%enditem

%beginitem
8. Що буде виведено за виконання фрагменту коду?

%code
print('R', end=' ')
j=9
while j>=1:
    j+=-2
print(j)
%code

\vspace{5mm}
%enditem

%beginitem
9. Що буде виведено за виконання фрагменту коду?
%code
if (19 <= 19)
    cout << "t";
else
    cout << "e";
%code

\vspace{5mm}
%enditem

%beginitem
10. Що буде виведено за виконання фрагменту коду?

%code
def g(a,b):
    c=45
    if a:
        c=1
    elif b<=-1:
         return 4
    else: 
        c=7
    return c

print(g(3,3))
%code
\vspace{5mm}
%enditem

%end
\newpage
\end{document}

